% =====================================================================
% 01_introduction.tex — introduction générale
% =====================================================================
\chapter{Introduction générale}
\label{chap:introduction}

\section{Contexte et motivation}

Les systèmes d'information bancaires figurent parmi les applications les plus
exigeantes en matière de modélisation, de contraintes d'intégrité et de
traçabilité. Une banque ne gère pas seulement des soldes : elle doit identifier
ses clients, suivre les conseillers et les agences, distinguer plusieurs
catégories de comptes, enregistrer les opérations, contrôler les droits des
signataires, appliquer des plafonds et des découverts, alerter sur les
échéances et conserver un historique des modifications sensibles. Le sujet du
module \module{} de la formation \formation{} demande précisément de concevoir
et d'implémenter un tel système sous Oracle, en exploitant les mécanismes du
modèle relationnel-objet \cite{sujetBDDA}.

Le projet ne se limite donc pas à une série de tables et de requêtes. Il doit
démontrer la maîtrise des types objet, de l'héritage, des références
\code{REF}, des collections \code{VARRAY} et \emph{nested tables}, des méthodes
membres et statiques, du polymorphisme, des types récursifs, des triggers et
des vues objet. Ces mécanismes sont imposés par le sujet et constituent le
cœur de l'évaluation. Le rapport présente la démarche suivie, les choix de
conception, les scripts réalisés, les résultats obtenus et les limites du
modèle.

\section{Problématique et objectifs}

La problématique générale peut être formulée ainsi : comment représenter, dans
une base Oracle, un domaine bancaire suffisamment riche pour illustrer les
apports du relationnel-objet, tout en garantissant les règles de gestion et la
cohérence des données ? Cette problématique se décompose en plusieurs
questions techniques :

\begin{itemize}
  \item comment modéliser une hiérarchie de comptes sans autoriser de compte
        générique, tout en conservant un traitement polymorphe des frais de
        tenue ;
  \item comment relier les objets par des références \code{REF} tout en
        maintenant l'intégrité référentielle, y compris pour un parrainage
        récursif ;
  \item comment stocker des collections de taille variable ou bornée
        (téléphones, signataires, statistiques mensuelles) et les interroger
        avec SQL ;
  \item comment centraliser les règles de sécurité dans des triggers sans
        provoquer d'erreur de table mutante ;
  \item comment exposer une vue objet qui agrège des conseillers et des
        comptes, et qui propose une méthode de mutation ;
  \item comment rendre l'ensemble rejouable, vérifiable et démontrable en
        soutenance.
\end{itemize}

Les objectifs du projet sont les suivants :

\begin{enumerate}
  \item produire un diagramme de classes UML complet et cohérent avec le
        sujet ;
  \item créer un utilisateur dédié, les types, les tables, les contraintes,
        les séquences et un jeu de données représentatif ;
  \item réaliser les manipulations de données demandées, notamment celles
        qui illustrent les contraintes différées et les collections ;
  \item écrire les requêtes SQL de la partie C, en particulier celles qui
        déréférencent des objets et interrogent des collections ;
  \item programmer les sous-programmes, méthodes et triggers de la partie D ;
  \item mettre en œuvre le type récursif, la collection multi-niveaux et la
        vue objet de la partie E ;
  \item valider le tout par un rejeu complet et documenter les résultats.
\end{enumerate}

\section{Périmètre fonctionnel et technique}

Le périmètre fonctionnel couvre les trois volets du sujet : les clients, les
conseillers et les agences ; les comptes et leurs opérations ; la
programmation et les fonctionnalités objet avancées. Les entités principales
sont les agences, les conseillers, les clients, les comptes courants, les
comptes d'épargne, les comptes à terme, les opérations, les alertes, l'audit,
les affectations et les relevés mensuels.

Le périmètre technique est volontairement limité à Oracle Database 21c
Express Edition, à la base enfichable \code{XEPDB1} et à l'utilisateur
\code{MASTER1}. Les scripts sont exécutés avec SQL*Plus. Aucune interface
graphique n'est développée : l'évaluation porte sur le schéma, le code SQL et
PL/SQL, les triggers et les objets. Cette limitation est assumée et discutée
dans les perspectives.

\begin{table}[htbp]
  \centering
  \small
  \caption{Environnement technique du projet.}
  \label{tab:environnement}
  \begin{tabularx}{\textwidth}{@{}L{4.2cm}X@{}}
    \toprule
    \entete{Élément} & \entete{Configuration retenue} \\
    \midrule
    SGBD & Oracle Database 21c Express Edition \\
    Base enfichable & \code{XEPDB1} \\
    Utilisateur applicatif & \code{MASTER1} \\
    Connexion & \code{sqlplus MASTER1/system@localhost:1521/XEPDB1} \\
    Langage procédural & PL/SQL \\
    Outil en ligne de commande & SQL*Plus \\
    Modélisation & UML, source PlantUML et Mermaid \\
    Versionnement & scripts SQL séparés par partie, plus un script unique \\
    \bottomrule
  \end{tabularx}
\end{table}

\section{Démarche méthodologique}

Le projet a été conduit selon une démarche en sept étapes, inspirée de la
méthode de conception d'une base de données. Cette démarche a guidé la
production des livrables et la vérification progressive des résultats.

\begin{table}[htbp]
  \centering
  \small
  \caption{Démarche de réalisation en sept étapes.}
  \label{tab:demarche}
  \begin{tabularx}{\textwidth}{@{}C{0.8cm}L{3.2cm}X@{}}
    \toprule
    \entete{\#} & \entete{Étape} & \entete{Contenu principal} \\
    \midrule
    0 & Cadrage & Compréhension du sujet, hypothèses, choix de conception et plan de travail. \\
    1 & Conception & Diagramme de classes UML complet, avec héritage, associations, cardinalités, collections et méthodes. \\
    2 & Schéma & Utilisateur, types, tables, contraintes, séquences et mécanismes de régénération. \\
    3 & Données & Jeu de test : agences, conseillers, clients, comptes, opérations et signataires. \\
    4 & Manipulation & Blocs B5a à B5f, procédures \code{change\_Portable} et \code{drop\_client}. \\
    5 & Requêtes & Requêtes 6.1 à 6.25, avec déréférencement, agrégats et collections. \\
    6 & Programmation & Sous-programmes catalogués, blocs anonymes, méthodes, triggers et transactions. \\
    7 & Intégration & Rejeu complet, guide de soutenance, préparation des captures et contrôle final. \\
    \bottomrule
  \end{tabularx}
\end{table}

Cette démarche présente l'avantage de séparer clairement la conception, la
création du schéma, l'exploitation et la validation. Elle permet aussi de
rejouer le projet partie par partie, ce qui est indispensable lors d'une
soutenance où le temps est limité.

\section{Organisation du rapport}

Le rapport est organisé en neuf chapitres, suivis d'une conclusion, de
perspectives, d'une bibliographie et de plusieurs annexes.

\begin{itemize}
  \item Le chapitre~\ref{chap:architecture} présente l'architecture générale
        et les choix de conception : partage entre objet et relationnel,
        héritage des comptes, références, collections, méthodes, triggers et
        vue objet.
  \item Le chapitre~\ref{chap:partieA} décrit la partie A : types métier,
        tables, contraintes, séquences, jeu de données et mécanismes de
        rejeu.
  \item Le chapitre~\ref{chap:partieB} présente les manipulations de la
        partie B, notamment la contrainte différée, la mise à jour d'un
        \code{VARRAY} et la suppression en cascade applicative.
  \item Le chapitre~\ref{chap:partieC} analyse les 25 requêtes SQL de la
        partie C et les techniques de déréférencement ou d'interrogation des
        collections.
  \item Le chapitre~\ref{chap:partieD1} traite des sous-programmes
        catalogués, des blocs anonymes et des méthodes de la partie D.
  \item Le chapitre~\ref{chap:partieD2} détaille les triggers, les
        transactions, la clôture de compte et les tests de sécurité.
  \item Le chapitre~\ref{chap:partieE} présente le type récursif, la
        collection multi-niveaux et la vue objet de la partie E.
  \item Le chapitre~\ref{chap:validation} rassemble la validation finale, la
        démonstration minutée, les 45 questions-réponses et les conseils de
        soutenance.
\end{itemize}

La conclusion générale dresse le bilan du projet. Les perspectives proposent
des extensions réalistes, clairement distinguées de ce qui a été réalisé. Les
annexes contiennent la liste des livrables, le plan détaillé des captures
d'écran, les hypothèses de conception et un glossaire.

\section{Conventions de lecture}

Pour faciliter la lecture, les noms de fichiers, de tables, de types et de
procédures sont composés en police à chasse fixe. Les extraits de code sont
numérotés et accompagnés d'une légende. Les références aux captures d'écran
prennent la forme \capture{05} : elles renvoient au plan de l'annexe~\ref{ann:captures},
qui précise la commande, le résultat attendu et l'emplacement prévu dans le
rapport.

Les résultats numériques cités dans le rapport correspondent aux sorties du
rejeu complet conservé dans \code{\_travail/\_FULL.out}. Lorsque la chronologie
des tests modifie l'état de la base, le rapport le signale explicitement. Cette
précaution est importante : la partie C s'exécute avant les tests destructifs
de la partie D, et certains comptes ou opérations ne sont donc plus visibles
après le rejeu complet.

\section*{Conclusion de l'introduction}
\addcontentsline{toc}{section}{Conclusion de l'introduction}

Le projet se situe à la rencontre de la modélisation UML, du SQL relationnel
et de la programmation PL/SQL orientée objet. Le rapport qui suit expose non
seulement le fonctionnement du système, mais aussi les raisonnements qui ont
conduit aux choix retenus. Cette explicitation est essentielle : le jury ne
doit pas seulement constater que le script s'exécute, il doit comprendre
pourquoi le schéma est conçu ainsi et ce que chaque partie démontre.
